% 2.2 — Les demandes de permission
\section{Les demandes de permission}
\label{sec:permissions}

La protection la plus répandue est la plus simple : avant une action sensible, l'agent
demande à la personne. Elle a deux faiblesses connues.

\paragraph{L'accoutumance.} Une demande répétée finit par être approuvée sans être lue.
Anthropic l'écrit à propos de Claude Code : avec le temps, les demandes mènent à une
« fatigue d'approbation, où les gens cessent de prêter attention à ce qu'ils
approuvent », et les utilisateurs approuvent 93\,\% des demandes de
permission~\cite{anthropic-automode}. Le phénomène est documenté bien avant les agents :
une étude de terrain sur plus de 25 millions d'avertissements de navigateur a mesuré que les
utilisateurs passaient outre l'avertissement SSL de Chrome dans 70,2\,\% des cas, et plus
vite pour les erreurs les plus fréquentes~\cite{akhawe-warningland}.

\paragraph{Confirmer n'est pas récupérer.} Une personne qui approuve \cmd{rm -f
Reports/report\_*.txt} en croyant que l'agent compte des fichiers
(section~\ref{sec:problematique}) n'a rien gagné à être consultée. Une fois la commande
passée, la demande n'offre aucun moyen de revenir en arrière.

\paragraph{Le juge automatique.} Pour réduire le nombre de demandes, Claude Code propose un
mode où un classifieur évalue chaque action avant son exécution et ne laisse passer sans
demande que ce qu'il juge aligné avec l'intention de l'utilisateur. Ses auteurs donnent
honnêtement ses limites : un taux de faux négatifs de 17\,\% sur des actions réellement
excessives, mesuré sur un petit échantillon, et le constat qu'il « ne remplace pas une
relecture humaine attentive sur une infrastructure à fort enjeu »~\cite{anthropic-automode}.

\begin{pourquoi}[Ce que \demo{} en retient]
Demander reste utile, mais \emph{seulement} quand la question est rare et qu'elle a un
sens : quand rien n'a pu être capturé et que la personne peut juger sur la commande
elle-même. Pour tout ce qui peut être copié avant, mieux vaut copier et ne pas déranger.
D'où deux choix de \demo{} : récupérer plutôt que demander quand c'est possible, et
distinguer, pour ce qui ne peut pas l'être, ce qu'on \emph{demande} de ce qu'on
\emph{refuse} (section~\ref{sec:ask-deny}).
\end{pourquoi}
